\documentclass[11pt]{article}
\usepackage[margin=23mm]{geometry}
\usepackage{amsmath,amssymb,amsthm}
\newcommand{\Q}{\mathbb Q}
\newcommand{\C}{\mathbb C}
\newcommand{\CH}{\mathrm{CH}}
\newcommand{\im}{\operatorname{im}}
\newtheorem{theorem}{Theorem}
\newtheorem{remark}[theorem]{Remark}
\title{Exact GST reduction: rational Chow-image equality to the complex Hodge-charge kernel}
\date{22 October 2026}
\begin{document}
\maketitle

\section{Input types, project-specific GST definitions}
Let X be an ACTUAL smooth projective complex variety of
dimension N, and p an integer with \(0\le p\le N\).
Set \(V_{\Q}=H^{2p}(X,\Q)\) and
\(V_{\C}=V_{\Q}\otimes_{\Q}\C\).
The projective rational canonical cycle-class map is
\[
 \operatorname{cl}^{can}_{X,p}:
       \CH^p(X)_{\Q}\longrightarrow V_{\Q}.
\]
The \texttt{GSTGradedWorldAlgebra.lean} laws set
\(D(a,b)=a+b\) and \(Q(a,b)=a-b\)
on the unbounded cosmic bidegree lattice.
The **separate, actual** geometric realization
of this charge is the Hodge-circle generator
\[
 \boxed{G_Xv=i(a-b)v\quad
            (v\in H^{a,b}(X)).}                   \tag{1}
\]
This complex-linear map is defined on \(V_{\C}\).
There is NO assumption that \(G_X(V_{\Q})\subset V_{\Q}\);
\(\ker G_X\cap V_{\Q}\) is, however, unambiguously typed.
\texttt{GSTClassicalHodgeFiberedCosmology.lean}
retains the genuine rational Hodge fiber and its
basis-index multiplicity, whereas its
\texttt{forgetMultiplicityToGST} forgets the latter.
The present theorem works on the entire genuine
bigraded Hodge realization, never the collapsed
rank-one abstract cosmic diagonal.

\section{Exactly the requested assumed rational statement}
Assume the input equation supplied for this reduction:
\[
 \boxed{
 V_{\Q}\cap\ker(G_X:V_{\C}\to V_{\C})
    =\im\bigl(\operatorname{cl}^{can}_{X,p}\bigr).
 }                                                 \tag{A}
\]
This is the rational Hodge conjecture's
algebraicity equality for this X,p,
expressed in Hodge-charge terminology.
It is a HYPOTHESIS in this research route,
not a consequence of charge conservation.

\begin{theorem}[All-N, all-p GST zero-charge kernel reduction]
For EVERY ACTUAL smooth projective complex X of
dimension N and every 0<=p<=N, the GST charge
realization (1) and the actual Hodge direct sum imply
\[
 \boxed{
 H^{p,p}(X)=
 \ker\!\left(G_X|_{H^{2p}(X,\C)}\right).
 }                                                \tag{B}
\]
Furthermore the universal spectral polynomial
\[
 \boxed{
 P_{X,p}:=\prod_{r=1}^{p}
   \left(I+\frac{G_X^2}{4r^2}\right)
 }                                                \tag{2}
\]
satisfies \(P_{X,p}^2=P_{X,p}\),
\(\im(P_{X,p})=H^{p,p}(X)\),
and \(P_{X,p}v=v\iff G_Xv=0\)
for every \(v\in V_{\C}\).
For p=0, the empty product is I.
\end{theorem}
\begin{proof}
Write the ACTUAL Hodge decomposition
\[
 V_{\C}=\bigoplus_{a+b=2p}H^{a,b}(X),
 \qquad v=\sum_{a+b=2p}v_{a,b}.
\]
By (1),
\[
 G_Xv=\sum_{a+b=2p}i(a-b)v_{a,b}.
\]
The terms lie in distinct summands of the
direct decomposition. Consequently \(G_Xv=0\)
iff every nonzero-charge term \(v_{a,b}\) vanishes.
Charge zero forces a=b, and a+b=2p forces a=b=p,
proving (B).

For a vector of charge \(q=a-b\), the scalar
eigenvalue of G is \(iq\).
In degree 2p, all possible nonzero charges are
\(\pm2,\pm4,\ldots,\pm2p\).
Thus the polynomial in (2) acts as one on charge
zero and zero on every allowed nonzero charge:
for \(q=\pm2s\) the factor with r=s vanishes.
Therefore \(P_{X,p}\) is the exact zero-charge
spectral projection, proving all assertions.
\end{proof}

\begin{theorem}[GST reduction of the INPUT Chow-image equation]
UNDER input (A), the exact combination with
(B) and the spectral projector yields
\[
 \boxed{
 \im(\operatorname{cl}^{can}_{X,p})
 =V_{\Q}\cap H^{p,p}(X)
 =\left\{\alpha\in V_{\Q}:P_{X,p}\alpha=\alpha\right\}.
 }                                                \tag{C}
\]
Equivalently, for EVERY rational \(\alpha\in V_{\Q}\),
\[
 \boxed{
 \alpha\text{ is canonically rational algebraic}
 \ \Longleftrightarrow\
 G_X\alpha=0
 \ \Longleftrightarrow\
 P_{X,p}\alpha=\alpha .
 }                                                \tag{3}
\]
\end{theorem}
\begin{proof}
The first equality of (C) is (A) with the
zero-charge kernel replaced by (B).
The second uses the spectral theorem above,
intersected with the ACTUAL rational lattice.
Thus every implication in (3) follows without
passing from rational vector spaces to an
unjustified complex span.
\end{proof}

\begin{remark}[Important logical direction and mathematical evidence]
(B) is PROVED from the GST charge realization and
Hodge decomposition, independently of assuming (A).
(A) then upgrades (B) to the rational
algebraicity criterion (C). One cannot infer
(A) merely by proving (B), since (B) does not
establish any CHOW representative existence.

Nor can one generally replace
\(V_{\Q}\cap H^{p,p}(X)\) by the entire complex
\(H^{p,p}(X)\) after tensoring with \(\C\).
For instance a projective K3 surface of
Picard rank one has
\(\dim_{\C}H^{1,1}=20\), whereas its rational
(1,1) subspace has dimension one.
The divisor case of the rational Hodge
conjecture holds, yet the complexification
of its rational Chow image has dimension
one, not twenty. This is a precise
coefficient-domain check, not a contradiction.
\end{remark}

\section{Executed symbolic checks and evidence status}
Wolfram exact evaluation tested every possible
charge \(q=-2p,-2p+2,\ldots,2p\) for
p=0,...,12. The eigenvalue polynomial gives
1 exactly at q=0 and 0 at the other charges,
for all such finite cases. The accompanying
complex-torus binomial Hodge multiplicities
were checked in 65 dimension/degree pairs
N=1,...,10, 0<=p<=N.
These finite checks validate implementation.
The complete all-N argument is the spectral
eigenvalue proof above. The reduction assumes
(A), rather than proving the Hodge conjecture.
\end{document}
